%------------------------------------------------------------------------------%

\usepackage{integracao_e_series}

\title{Apresentação da Disciplina}

%------------------------------------------------------------------------------%
\begin{document}

\addtitle

%------------------------------------------------------------------------------%
\section{Introdução}

%------------------------------------------------------------------------------%
\begin{frame}{Disciplina}
   Disciplina: \emph{\large Integrais e Séries}\\[7mm]
   Professor: \emph{\large Luis A. D'Afonseca}
\end{frame}

%------------------------------------------------------------------------------%
\begin{frame}{Disciplina}
\begin{itemize}[<+->]
  \item Aula semi invertida
  \item Estude o conteúdo com antecedência
  \item As apresentações são um resumo da teoria -- \emph{Estude o livro}
  \item Traga papel e lápis para fazer os exercícios em aula
  \item Estude o conteúdo seriamente entre uma aula e outra
  \item Não deixe o conteúdo acumular
  \item Não maratone o estudo de Matemática
  \item Preencha as lacunas do seu conhecimento de Matemática \\
        \href{https://pt.khanacademy.org}{Khan Academy}
\end{itemize}
\end{frame}

%------------------------------------------------------------------------------%
\section{Integração}

\begin{frame}{Integração}
  \begin{enumerate}[<+->]
    \item Cálculo da área de regiões curvas
    \begin{itemize}
      \item[] Integrais Definidas
    \end{itemize}
    \item Operação inversa a derivação
    \begin{itemize}
      \item[] Integrais Indefinidas
    \end{itemize}
    \item Teorema Fundamental do Cálculo
    \item Integrais Impróprias
    \item Cálculo de Áreas e Volumes \\[5mm]
    \item Técnicas de Integração
  \end{enumerate}
\end{frame}

%------------------------------------------------------------------------------%
\begin{frame}{Prerrequisitos}
  \begin{itemize}[<+->]
    \item Manipulações algébricas
    \item Polinômios
    \item Funções
    \item Limites
    \item Derivadas
  \end{itemize}
\end{frame}

%------------------------------------------------------------------------------%
\section{Séries}

\begin{frame}{Séries}
  \begin{enumerate}[<+->]
    \item Sequências numéricas \tabto{40mm} -- lista infinita de números
    \item Séries numéricas     \tabto{40mm} -- Soma dos termos de uma sequência \\[5mm]
    \item Séries de potências  \tabto{40mm} -- ``Polinômio de grau infinito''
    \item Séries de Taylor     \tabto{40mm} -- Como representar funções como séries de potências
  \end{enumerate}
\end{frame}

%------------------------------------------------------------------------------%
\begin{frame}{Prerrequisitos}
  Além dos prerrequisitos para integração
  \begin{itemize}[<+->]
    \item Integração
    \item Fatorial
    \item Módulo
    \item Notação de somatório
  \end{itemize}
\end{frame}

%------------------------------------------------------------------------------%
\section{Material Didático}

%------------------------------------------------------------------------------%
\begin{frame}{Material Didático}
  \begin{enumerate}[<+->]
    \item Apostila (rascunho)  ``\href{https://sites.google.com/view/prof-luis-dafonseca/gradua\%C3\%A7\%C3\%A3o/Integrais-e-Series}{Integração e Séries}''
    \item \href{https://drive.google.com/drive/u/2/folders/1Bb1hqfdwt6E3jZpB8L4HRhbtW-UopeNB}{Material complementar}
    \item Thomas ou Stewart
    \item Material online
  \end{enumerate}
\end{frame}

%------------------------------------------------------------------------------%
\section{Lista Mínima}

%------------------------------------------------------------------------------%
\begin{frame}{Lista Mínima}
  Estudar os prerrequisitos da disciplina

  Estudar o Apêndice A da Apostila

  Ler o Apêndice B da Apostila

  \vfill
  Atenção: A prova é baseada no livro, não nas apresentações
\end{frame}

%------------------------------------------------------------------------------%
\end{document}
%------------------------------------------------------------------------------%
